% =======================================================
%
%(iv) Synergistic Activities
%Each individual identified as a senior/key person must provide a document of up to one-page that includes a list of up to five distinct examples that demonstrates the broader impact of the individual's professional and scholarly activities that focus on the integration and transfer of knowledge as well as its creation. Senior/key personnel must prepare, save, and submit these documents as part of their proposal via Research.gov or Grants.gov. Examples may include, among others: innovations in teaching and training; contributions to the science of learning; development and/or refinement of research tools; computation methodologies and algorithms for problem-solving; development of databases to support research and education; broadening the participation of groups underrepresented in STEM; participation in international research collaborations; participation in national and/or international standards development efforts; and service to the scientific and engineering community outside of the individual's immediate organization.
% 
% - Complete one document for all investigators/senior personnel working on the project.
% - Each plan must be no more than 1 page.
% - Provide up to 5 examples of synergistic activities.
% - Assure each example of synergistic activity is distinct from other examples. 
% - The file must be converted into PDF prior to upload to Research.gov
%
% For proposals submitted via Research.gov, the system will automatically paginate a proposal. Each section of the proposal that is uploaded as a file should leave out page numbering unless otherwise directed within Research.gov.
%
% Allowed fonts:
% - Arial[7] (not Arial Narrow), Courier New, or Palatino Linotype at a font size of 10 points or larger;
% - Times New Roman at a font size of 11 points or larger; or
% - Computer Modern family of fonts at a font size of 11 points or larger.
% 
% =======================================================

\documentclass{NSF}
\begin{document}

\clearpage
\setcounter{page}{1}

% ======================================
\section*{Appendix 5: Synergistic Activities \hfill John A.\ Marohn}
% ======================================

\begin{enumerate}

\item {\bf Technology Transfer}.
\emph{Atomic force microscopy apparatus, methods, and applications}, John A. Marohn, Sarah R. Nathan, and Ryan P. Dwyer. U.S. Patent 11,808,783 issued November 7, 2023.
This patent covers Marohn and co-workers' invention of the ``phase-kick electric force microscope''.
The patent describes an atomic force microscopy (AFM) apparatus and method that enable observing photoconductivity transients with nanometer spatial resolution and nanosecond to picosecond time resolution, the timescale relevant for studying the recombination of photo-generated charges in the world's highest efficiency photovoltaic films.
The apparatus includes an AFM, a light source for illumination of a sample operatively coupled to the AFM, a voltage source operatively coupled to the AFM, and a control circuitry operatively coupled to the light source and the voltage source. 
The apparatus improves the time resolution and acquisition time of photoconductivity transients and is applicable to a wide array of energy-harvesting and energy-storage materials.

\item {\bf Open Source Software}.
Marohn and his research team have released many software packages to GitHub.
A representative example is the Python package \emph{dissipationtheory}, \url{https://github.com/Marohn-Group/dissipationtheory}, consisting of over 19,000 lines of Python code and 75 Jupyter notebooks of example calculations. 
The package computes the frequency shift and dissipation experienced by a charged scanning probe microscope tip over a metal, dielectric, or semiconductor sample. 
The calculation uses a point-probe, reaction-field model of the tip-sample interaction, introduced by Lekalla and Loring, and an image-charge extension of this model, invented by Marohn and Loring, that is applicable to a probe represented as a cylindrically symmetric equipotential surface.
The interaction of the sample with the cantilever can thus be modeled analytically, approximating the cantilever tip as a point charge, or numerically, approximating the cantilever tip as a sphere or a sphere plus a cone.
The cantilever can oscillate either parallel or perpendicular to the sample surface, and the cantilever voltage can be static or oscillatory.
This code is released under the MIT License.

\item {\bf Innovations in Teaching}.
To facilitate the teaching of graduate quantum mechanics, Marohn created the Mathematica package \emph{UniDyn}, \url{https://github.com/JohnMarohn/UniDyn}.  
This package introduces an algorithm for symbolically computing the unitary evolution of quantum mechanical operators, with applications in magnetic resonance, quantum optics, quantum computing, and electron transfer theory. 
The package includes example analytical calculations describing
spin magnetization evolving under off-resonance, variable-phase rf irradiation; 
the free-induction decay magnetic resonance signal for two weakly coupled spins; 
magnetic resonance signal for two weakly coupled spins in the INEPT polarization transfer, heteronuclear COSY, and homonuclear COSY experiments; 
harmonic oscillator position and momentum evolution under force, free evolution, and squeezing Hamiltonians; 
unitary rotations relevant to quantum optics; and 
unitary rotations relevant to Marcus electron-transfer theory.
This code is released under the GNU General Public License, Version 3.

\item {\bf International Involvement}.
Marohn created, secured NSF and Kavli Foundation funding for, and chaired the first and second International Summer Schools on Nanoscale Magnetic Resonance Imaging (2006 and 2009; Ithaca, New York).
Subsequent ``nano-MRI'' conferences were hosted by teams in France (2010), Switzerland (2012), Canada (2015), France (2018), and Spain (2021).
The seventh conference will be held in the Netherlands in 2026.

\item {\bf Mentoring}.
Since 1999, Marohn has advised 10 postdocs (1 current), 38 Ph.D.\ students (7 current),  6 M.S.\ students, and 29 undergraduate researchers.
Awards garnered by Marohn's graduate students and postdocs include
the Congressional Science and Engineering Fellowship of the Materials Research Society and the Optical Society of America; 
a DOE Sunshot Postdoctoral Fellowship; 
two NSF Graduate Research Fellowship awards; and
chairing the Gordon Research Seminar on Unconventional Semiconductors and Their Applications.

\end{enumerate}


\end{document}